\begin{frame}{Phân tích điểm chức năng (FPA)}
    \begin{block}{}
        \begin{itemize}
            \item Phân tích điểm chức năng (FPA) được phát triển
            bởi Allan Albrecht vào cuối những năm 1970 tại IBM.
            \item International Function Point Users Group (IFPUG).
            \item Được quy định bởi tiêu chuẩn ISO 20296, ISO/IEC 14143.
            \item Chia nhỏ yêu cầu thành bốn loại di chuyển dữ liệu: Entry (E), Exit (X), Read (R), Write (W).
        \end{itemize}
    \end{block}
    \Large Dùng để ước lượng khối lượng công việc (Function Points).
\end{frame}

%% TODO use case points

%% Traceability Matrix
\begin{frame}{Ma trận theo dõi (Traceability Matrix)}
    \begin{block}{}
        \begin{quote}
            Ma trận theo dõi là một công cụ quản lý dự án
            để theo dõi quan hệ giữa các yêu cầu và các yếu tố
            thiết kế của hệ thống.\\
            %% Use Cases, Non-Functional Requirements, Test Cases, Classes, Packages, Servers, Containers, Components, Nodes, etc.
            \small Ca sử dụng, yêu cầu phi chức năng, ca kiểm thử, lớp, gói, máy chủ, container, thành phần, node, v.v.
        \end{quote}
    \end{block}
    \begin{figure}
        \centering
        \includegraphics[width=0.5\textwidth]{matrix}
    \end{figure}
\end{frame}

\begin{frame}{Kiểm tra (verification) and Thẩm định (validation)}
    \begin{block}{Kiểm tra (Do it right?)}
        Kiểm tra là quá trình đảm bảo rằng phần mềm được
        xây dựng đúng theo các yêu cầu và đặc tả kỹ thuật.
    \end{block}
    \begin{block}{Thẩm định (Do the right thing?)}
        Thẩm định là quá trình đảm bảo rằng các yêu cầu và
        đặc tả kỹ thuật của phần mềm đáp ứng đúng nhu cầu và
        mong đợi của người dùng.
    \end{block}
\end{frame}

\section{Yêu cầu phi chức năng (Non-Functional)}

\begin{frame}{Yêu cầu phi chức năng (Non-Functional Requirements)}
    \begin{block}{}
        \begin{itemize}
            \item Hiệu suất (Performance, Capacity)
            \item Sẵn sàng (Availability)
            \item Bảo mật (Security)
            \item Khả năng phục hồi (Recoverability)
            \item Dễ sử dụng (Usability)
            \item Bảo trì (Maintainability)
            \item Di động (Mobile)
            \item Đáng tin cậy (Reliability)
            \item Tích hợp (Integration)
            \item V.v.
        \end{itemize}
    \end{block}
\end{frame}

\begin{frame}{Tính sẵn sàng}
    \begin{block}{}
        \begin{itemize}
            \item Sẵn sàng (Availability) là tỷ lệ thời gian hệ thống
            hoạt động so với tổng thời gian.
            \[A = \dfrac{E_\text{up}}{E_\text{down} + E_\text{up}}\]
        \end{itemize}
    \end{block}
\end{frame}

%% capacity
\begin{frame}{Hiệu suất}
    \begin{block}{}
        \begin{itemize}
            \item Hiệu suất (Capacity) là số lần giao tiếp mà hệ thống
            có thể xử lý trong một đơn vị thời gian.
            \[CPS = \text{Clicks Per Second}\]
        \end{itemize}
    \end{block}
\end{frame}

%% recovery time objective
\begin{frame}{Khả năng phục hồi}
    \begin{block}{}
        \begin{itemize}
            \item Khả năng phục hồi (Recovery) là thời gian cần thiết
            để hệ thống phục hồi sau một sự cố.
            \[RTO = \text{Recovery Time Objective}\]
        \end{itemize}
    \end{block}
\end{frame}

%% maintainability
\begin{frame}{Bảo trì}
    \begin{block}{}
        \begin{itemize}
            \item Bảo trì (Maintainability) là thời gian cần thiết
            để sửa một lỗi trong hệ thống.
            \[MTTR = \text{Mean Time To Repair}\]
        \end{itemize}
    \end{block}
\end{frame}

%% good and bad NFRs
\begin{frame}{Yêu cầu phi chức năng}
    \begin{columns}
    \column{0.5\textwidth}
        {\Huge \color{red}Không tốt}
        \begin{itemize}
            \item Phần mềm phải nhanh.
            \item Phải dễ bảo trì.
            \item Giao diện người dùng phải hấp dẫn.
        \end{itemize}
    \column{0.5\textwidth}
        {\Huge \color{blue}Tốt}
        \begin{itemize}
            \item Được cài đặt trên máy chủ kiểm thử (xem Phụ lục A),
            phần mềm phải phản hồi trong vòng 20ms
            cho các yêu cầu trong UC1--UC7.
            \item Thời gian tối đa để sửa một lỗi
            không quá hai giờ.
            \item Ít nhất 80\% người dùng beta phải xác nhận
            ẩn danh rằng giao diện người dùng đủ hấp dẫn.
        \end{itemize}
    \end{columns}
\end{frame}

%% các lỗi chính trong mô tả yêu cầu
\begin{frame}{Các lỗi chính trong mô tả yêu cầu}
    \begin{block}{}
        \begin{itemize}
            \item Không có từ điển thuật ngữ hoặc viết từ điển tồi.
            \item Dùng câu hỏi, thảo luận, ý kiến.
            \item Lẫn lộn yêu cầu chức năng và phi chức năng.
            \item Lẫn lộn yêu cầu và thông tin bổ sung.
            \item Yêu cầu phi chức năng không đo được.
            \item Đưa kèm hướng dẫn cài đặt, triển khai.
            \item Thiếu góc nhìn của người dùng.
            \item Nhiễu.
            \item Không rõ ràng (Who does What).
        \end{itemize}
    \end{block}
\end{frame}

\section{Ước lượng}

\begin{frame}{Ước lượng}
    \begin{columns}
        \column{0.3\textwidth}
        \begin{figure}
            \centering
            \includegraphics[width=0.8\textwidth]{../images/books/boehm}
        \end{figure}
        \column{0.7\textwidth}
            \only<1>{
                \begin{block}{}
                    Trong các quyết định thiết kế,
                    yếu tố quan trọng và khó
                    khăn nhất là ước tính
                    chi phí của một dự án phần mềm.\\
                    \textbf{Barry W. Boehm}, \emph{Software Engineering Economics}
                \end{block}
            }
    \end{columns}
\end{frame}

\begin{frame}{Ước lượng}
    \begin{block}{}
        \begin{enumerate}
            \item Ước lượng kích thước theo kilo-line-of-code
            \[K = FP \times \text{hệ số chuyển đổi}\]
            \item Tìm hệ số điều chỉnh năng lực ($F$).
            \item Tìm các hệ số $a$, $b$, $c$ và $d$ từ mô hình COCOMO II.
            \item Tính kích thước theo người-tháng (MM - man-month):
            \[E = a \times K^b \times F\]
            \item Tính thời gian theo tháng (D):
            \[D = c \times E^d\]
        \end{enumerate}
    \end{block}
\end{frame}

%% hệ số chuyển đổi
\begin{frame}{Hệ số chuyển đổi}
    \begin{block}{}
        \begin{itemize}
            \item Hệ số chuyển đổi (Conversion Factor) là tỷ lệ
            giữa kích thước ước lượng và kích thước thực tế.
            \[KLOC = FP \times \text{hệ số chuyển đổi}\]
            \item Hệ số chuyển đổi theo ngôn ngữ (KLOC/FP)
            \begin{itemize}
                \item Java: 0.4
                \item C\#: 0.5
                \item C: 0.6
                \item Python: 0.2
                \item COBOL: 0.8
            \end{itemize}
        \end{itemize}
    \end{block}
\end{frame}

%% hệ số điều chỉnh năng lực
\begin{frame}{Hệ số điều chỉnh năng lực}
    \begin{block}{}
        \begin{itemize}
            \item Hệ số điều chỉnh năng lực (Effort Adjustment Factor) là
            hệ số điều chỉnh dựa trên một số yếu tố ảnh hưởng
            đến khả năng thực hiện của dự án.
            \[F = \prod_{i=1}^n \text{hệ số hiệu chỉnh của yếu tố}_i\]
            \item Mô hình COCOMO II có 17 yếu tố ảnh hưởng.
        \end{itemize}
    \end{block}
\end{frame}

%% các hệ số a, b, c, d
\begin{frame}{Các hệ số $a$, $b$, $c$, $d$}
    \begin{block}{}
        \begin{itemize}
            \item Các hệ số $a$, $b$, $c$, $d$ được xác định từ
            mô hình COCOMO II.
            \item Các hệ số này phụ thuộc vào các pha của dự án
            (Early Design, Post-Architecture, Reuse Model).
            \item Các hệ số này được xác định từ bảng.
        \end{itemize}
    \end{block}
\end{frame}

%% cone of uncertainty
\begin{frame}{Nón không chắc chắn (Cone of Uncertainty)}
    \begin{figure}
        \centering
        \includegraphics[width=0.6\textwidth]{cone}
    \end{figure}
    \Large\color{red} "Phần mềm này sẽ mất bao nhiêu tiền?"\\
    "Việc phát triển sẽ kéo dài và tiêu hết tiền của bạn".
\end{frame}

%% books
\begin{frame}{Tài liệu tham khảo}
    \begin{columns}
        \column{0.5\textwidth}
        \begin{figure}
            \centering
            \includegraphics[width=0.4\textwidth]{../images/books/wiegers}
            \caption{Karl Wiegers, \emph{Software Requirements}}
        \end{figure}
        \column{0.5\textwidth}
        \begin{figure}
            \centering
            \includegraphics[width=0.4\textwidth]{../images/books/wiegers2}
            \caption{Karl Wiegers, \emph{More About Software Requirements}}
        \end{figure}
    \end{columns}
\end{frame}

\begin{frame}{Tài liệu tham khảo}
    \begin{columns}
        \column{0.5\textwidth}
        \begin{figure}
            \centering
            \includegraphics[width=0.4\textwidth]{../images/books/cockburn}
            \caption{Alistair Cockburn, \emph{Writing Effective Use Cases}}
        \end{figure}
        \column{0.5\textwidth}
        \begin{figure}
            \centering
            \includegraphics[width=0.4\textwidth]{../images/books/mcconnell}
            \caption{Steve McConnell, \emph{Software Estimation: Demystifying the Black Art}}
        \end{figure}
    \end{columns}
\end{frame}
\documentclass{beamer}

% Vietnamese support
\usepackage{fontenc}
\usepackage[utf8]{vietnam}

\newcommand{\ul}[1]{\emph{#1}}

% Beamer theme
\usetheme{Berlin}

\title{Thực hành phát triển hệ thống Trí tuệ nhân tạo\\
RUP, Agile/XP
}
\author{Trần Quốc Long}
\institute{Trường ĐH Công nghệ, ĐHQG Hà Nội}
\date{\today}

\begin{document}

\begin{frame}
    \titlepage
\end{frame}

\begin{frame}{Mục lục}
    \tableofcontents
\end{frame}

\section{Các phương pháp phát triển phần mềm}

\begin{frame}{Các phương pháp phát triển}
    \begin{itemize}
        \item Quy trình thác nước (Waterfall/SDLC)
        \item Quy trình phát triển thống nhất (RUP)
        \item Khung phần mềm của Microsoft (MSF)
        \item Thư viện cơ sở hạ tầng CNTT (ITIL)
        \item Phong cách Agile và các phương pháp khác
    \end{itemize}
\end{frame}

\begin{frame}{Quy trình thác nước}
    \begin{figure}
        \centering
        \includegraphics[width=0.8\textwidth]{waterfall}
    \end{figure}
\end{frame}

\begin{frame}{Vòng đời phần mềm (SDLC)}
    \begin{figure}
        \centering
        \includegraphics[width=0.8\textwidth]{sdlc}
    \end{figure}
\end{frame}

\section{RUP}

\begin{frame}{Quy trình phát triển thống nhất (RUP)}
    Lặp và tăng trưởng
    \begin{figure}
        \centering
        \includegraphics[width=0.8\textwidth]{rup}
    \end{figure}
\end{frame}

\begin{frame}{Lặp (Iterative)}
    \begin{figure}
        \centering
        \includegraphics[width=0.8\textwidth]{iterative}
    \end{figure}
\end{frame}

\begin{frame}{Tăng trưởng (Incremental)}
    \begin{figure}
        \centering
        \includegraphics[width=0.8\textwidth]{i-and-i}
    \end{figure}
\end{frame}

\begin{frame}{Vai trò trong RUP}
    \begin{figure}
        \centering
        \includegraphics[width=0.7\textwidth]{rup-architect}
    \end{figure}
    \begin{figure}
        \centering
        \includegraphics[width=0.7\textwidth]{rup-designer}
    \end{figure}
\end{frame}

%% Microsoft
\begin{frame}{Khung phần mềm của Microsoft (MSF)}
    \begin{figure}
        \centering
        \includegraphics[width=0.9\textwidth]{msf}
    \end{figure}
\end{frame}

\begin{frame}{Thư viện cơ sở hạ tầng CNTT (ITIL)}
    \begin{figure}
        \centering
        \includegraphics[width=0.5\textwidth]{itil}
    \end{figure}
\end{frame}

\begin{frame}{Các phương pháp khác}
    \begin{columns}
        \column{0.5\textwidth}
        \begin{figure}
            \centering
            \includegraphics[width=0.7\textwidth]{spiral}
            \caption{Spiral}
        \end{figure}
        \column{0.5\textwidth}
        \begin{figure}
            \centering
            \includegraphics[width=0.5\textwidth]{vmodel}
            \caption{V-Model}
        \end{figure}
        \begin{figure}
            \centering
            \includegraphics[width=0.5\textwidth]{rad}
            \caption{Rapid Application Development}
        \end{figure}
    \end{columns}
\end{frame}

\section{Agile}

\begin{frame}{Agile}
    \begin{columns}
        \column{0.3\textwidth}
        \begin{figure}
            \centering
            
\includegraphics[width=0.8\textwidth]{../images/books/agile-martin}
        \end{figure}
        \column{0.7\textwidth}
            \only<1>{
                \begin{block}{}
                    Trong Agile, bức tranh tổng thể phát triển
                    cùng với phần mềm. Với mỗi lần lặp,
                    thiết kế của hệ thống được cải thiện 
                    tốt nhất có thể tại thời điểm đó.
                    Không cần dành nhiều thời gian để nhìn trước
                    các yêu cầu và nhu cầu trong tương lai.\\
                    \textbf{Robert Martin}, \emph{Agile Software Development. Principles, Patterns, and Practices}
                \end{block}
            }
    \end{columns}
\end{frame}

\begin{frame}{Quy trình Agile}
    \begin{figure}
        \centering
        \includegraphics[width=0.7\textwidth]{agile}
    \end{figure}
\end{frame}

\begin{frame}{Công cụ Kanban}
    \begin{figure}
        \centering
        \includegraphics[width=0.9\textwidth]{kanban}
    \end{figure}
\end{frame}

\begin{frame}{Khung Scrum}
    \begin{figure}
        \centering
        \includegraphics[width=\textwidth]{scrum}
    \end{figure}
\end{frame}

\begin{frame}{Lập trình cực đoan (XP)}
    \begin{columns}
        \column{0.3\textwidth}
        \begin{figure}
            \centering
            \includegraphics[width=0.8\textwidth]{../images/books/xp}
        \end{figure}
        \column{0.7\textwidth}
            \only<1>{
                \begin{block}{}
                    Liên tục tinh chỉnh thiết kế
                    của hệ thống, bắt đầu từ một khởi đầu rất đơn giản.
                    Loại bỏ bất kỳ sự linh hoạt nào không hiệu quả.\\
                    \textbf{Kent Beck}, \emph{Extreme Programming Explained: Embrace Change}
                \end{block}
            }
    \end{columns}
\end{frame}

\begin{frame}{Lập trình cực đoan (XP)}
    \begin{figure}
        \centering
        \includegraphics[width=0.65\textwidth]{xp}
    \end{figure}
\end{frame}

\begin{frame}{Chi phí của bug}
    \begin{figure}
        \centering
        \includegraphics[width=0.9\textwidth]{cost-of-bug}
    \end{figure}
\end{frame}

\section{Các yếu tố chất lượng của thiết kế}

\begin{frame}{Tính đơn giản}
    \begin{itemize}
        \item \ul{\color{red}Đầu tư} nhỏ nhất có thể vào thiết kế trước khi
        nhận được lợi nhuận từ nó. -- \textbf{Kent Beck}
        \item Thiết kế chứa \ul{\color{red}sự phức tạp không cần thiết} khi
        có các yếu tố không hữu ích. Điều này thường xảy ra khi
        ta dự đoán các thay đổi có thể có đối với yêu cầu và tìm cách
        xử lý những thay đổi tiềm năng đó. -- \textbf{Robert Martin}        
        \item Không làm chủ sự phức tạp của phần mềm 
        có hậu quả là các dự án bị trễ, vượt ngân sách và thiếu sót
        so với yêu cầu. Tình trạng này gọi là
        \ul{\color{red}khủng hoảng phần mềm}, nhưng thẳng thắn mà nói,
        một căn bệnh kéo dài như vậy phải được gọi là
        trạng thái bình thường. -- \textbf{Grady Booch}
    \end{itemize}
\end{frame}

\begin{frame}{Tính đơn giản}
    \begin{figure}
        \centering
        \includegraphics[width=0.9\textwidth]{complexity1}
    \end{figure}
\end{frame}

\begin{frame}{Tính đơn giản}
    \begin{figure}
        \centering
        \includegraphics[width=0.8\textwidth]{complexity2}
    \end{figure}
\end{frame}

\begin{frame}{Tính đơn giản}
    \begin{block}{}
        Đức tính chính của một kiến trúc sư là khả năng
        \ul{giảm bớt sự phức tạp}. Một kiến trúc sư giỏi
        không bao giờ tự hào về một sơ đồ phức tạp. Thay vào
        đó, họ sẽ tự hào về một bản vẽ đơn giản và dễ hiểu
        với vài hình chữ nhật mà giải thích hoàn hảo một
        ứng dụng đa tầng. Việc này thực sự khó. Đó
        là nơi mà tư duy kiến trúc thực sự tỏa sáng
        -- \textbf{Yegor Bugayenko}
    \end{block}
\end{frame}

\begin{frame}{Tính không trùng lặp}
    \begin{itemize}
        \item Khi có \ul{\color{red}mã dư thừa} trong hệ thống, việc thay
        đổi hệ thống trở nên khó khăn. Các lỗi được tìm thấy
        phải được sửa trong mọi lần lặp lại. Tuy nhiên,
        vì mỗi lần lặp lại hơi khác nhau, nên việc sửa
        lỗi không phải lúc nào cũng giống nhau -- \textbf{Robert Martin}
    \end{itemize}
\end{frame}

\begin{frame}{Tính mô-đun hóa}
    \begin{itemize}
        \item Một thiết kế là \ul{\color{red}bất động} nếu nó
        chứa các thành phần
        được dùng trong các hệ thống khác, nhưng nỗ lực và
        rủi ro để tách các phần đó ra khỏi hệ thống ban đầu
        là quá lớn -- \textbf{Robert Martin}
        \item Một thiết kế là \ul{\color{red}cứng} nếu một
        thay đổi đơn lẻ gây ra một loạt các thay đổi tiếp
        theo trong các module phụ thuộc.
        Càng nhiều module phải thay đổi,
        thiết kế càng \ul{cứng} -- \textbf{Robert Martin}
    \end{itemize}
\end{frame}

\begin{frame}{Tài liệu tham khảo}
    \begin{columns}
        \column{0.5\textwidth}
        \begin{figure}
            \centering
            \includegraphics[width=0.4\textwidth]{../images/books/rup}
            \caption{Per Kroll và cộng sự, \emph{A Practitioner's Guide to the RUP}}
        \end{figure}
        \column{0.5\textwidth}
        \begin{figure}
            \centering
            
\includegraphics[width=0.4\textwidth]{../images/books/agile-martin}
            \caption{Robert Martin,
            \emph{Agile Software Development. Principles, Patterns, and Practices}}
        \end{figure}
    \end{columns}
\end{frame}

\begin{frame}{Tài liệu tham khảo}
    \begin{columns}
        \column{0.5\textwidth}
        \begin{figure}
            \centering
            \includegraphics[width=0.4\textwidth]{../images/books/clean-architecture}
            \caption{Robert Martin, \emph{Clean Architecture: A Craftsman's Guide to Software Structure and Design}}
        \end{figure}
        \column{0.5\textwidth}
        \begin{figure}
            \centering
            \includegraphics[width=0.4\textwidth]{../images/books/code-complete}
            \caption{Steve McConnell, \emph{Code Complete}}
        \end{figure}
    \end{columns}
\end{frame}

\begin{frame}{Tài liệu tham khảo}
    \begin{columns}
        \column{0.5\textwidth}
        \begin{figure}
            \centering
            \includegraphics[width=0.4\textwidth]{../images/books/brooks}
            \caption{Frederick Brooks Jr., \emph{Mythical Man-Month, The: Essays on Software Engineering}}
        \end{figure}
        \column{0.5\textwidth}
        \begin{figure}
            \centering
            \includegraphics[width=0.4\textwidth]{../images/books/pragmatic}
            \caption{David Thomas và cộng sự, \emph{The Pragmatic Programmer: Your Journey To Mastery}}
        \end{figure}
    \end{columns}
\end{frame}

\end{document}
